% Independent collation and translator's mathematical checks; not author text.
\begingroup
\section*{D19-R：第9、10节的法英对读与核查}
\markboth{第十九版：第9、10节回校}{第十九版：第9、10节回校}
\addcontentsline{toc}{section}{D19-R：第9、10节的法英对读与核查}

\noindent\textbf{核查范围。}逐段对读法稿物理53--60页上部、英译36--41页第11节以前；新正文为 \texttt{v19\_section9.tex} 和 \texttt{v19\_section10.tex}。原第十八版模块仍完整保留。本报告中的补证及反例均为译校工作，不称为恢复的作者证明。法稿的手写页码分别比PDF物理页码小3；不把法稿第52页误当PDF物理52页。

\subsection*{一、可核对的异文及归责}
\begin{longtable}{p{.14\textwidth}p{.19\textwidth}p{.55\textwidth}}
\toprule
位置 & 法稿物理页 & 本轮依据与处理\\\midrule\endhead
9.1 & 53 & $E$、$\mathcal O_P(n)$ 可辨；关系符靠扫描边界，仍不宣称逐笔恢复。典范同构有独立依据。二次曲面的旧维数条件被划去，后改为 $P^3$ 中的例子。\\
9.2 & 53--54 & 参数空间在法稿上有对偶记号；漏掉它的是英译。8.8缺完整假设的问题未解决。\\
9.3(c) & 54 & $n\ge3$ 后的手写“ou”可辨，恢复为“或”；共同条件 $n\ge2$ 不变。\\
9.3证明 & 54--55 & 法稿为 $f(x)$ 的仿射邻域、小写 $k(t_1,\ldots,t_r)$ 和 $f(X)\not\subset L$。英译的三处写法须据此订正。下标 $i/s$ 混用则两本都有。\\
9.3末 & 55 & 划去的“切映射为单射”论证与最后一条 $n=2$ 备注分别保存删除身份。英译38页的三个平方计算保留并明确其来源。\\
9.4--9.5 & 55 & 9.4只有编排提示；9.5仍无完整定理。可辨的几何整条件、大写 $X\to P$ 和奇异点簇恢复；中间未辨词不猜补。\\
10.1 & 56 & 法稿两处有 $\overline Z$，不可继续一概归为漏上划线。单条直线判据仍须明确闭包；该句末法稿误打为 $k$，英译作 $Z$ 并加P.B.批注。\\
10.2 & 57--58 & $T=X\times_P L^0$ 及到 $J/J^2$ 的完整同态据法稿恢复。“主要用于”订正旧译“只用于”。奇异点在 $k$ 上有理的过强措辞两本共有。\\
10.3 & 59 & 爆破记号 $\widetilde X$、大写 $P$、“零点预概形”均可从法稿恢复；不再列为作者混记或缺词。整条仍为已划去的草稿。\\
10.4 & 59--60 & $\dim X\le1$、几何纤维的 $\operatorname{Spec}k'$、逐纤维“在 $s$ 上方”均原本存在，英译分别误作 $\ge1$、$\operatorname{Spec}k$ 及漏去范围。\\
\bottomrule
\end{longtable}

\subsection*{二、9.3原条件不足：爆破的固定分支反例}
取代数闭域 $k$，$p=[1:0:0]\in\mathbf P^2_k$，令
\[
 f:X=\operatorname{Bl}_p\mathbf P^2_k\longrightarrow\mathbf P^2_k,
 \qquad E=f^{-1}(p)\simeq\mathbf P^1_k.
\]
$X$ 光滑、射影且几何整，像维数为2。任取 $x\in E(k)$。所有经过 $p$ 的 $n$ 次曲线方程都属于 $I_p$，故拉回后都含 $E$。泛方程不是零式，也不只支撑在 $p$；它的严格变换是另一个非空分支，因此总变换可约。原9.3(a)已经失败。

取 $p$ 的仿射邻域，坐标记作 $(a,b)$；爆破的一张图为
\[
 a=u,\qquad b=uv,\qquad E=V(u),\qquad x=(0,0).
\]
例如 $a+a^2+b^2$ 的拉回为
\[
 u\bigl(1+u+uv^2\bigr).
\]
令(b)的线性子簇为 $L=V(b)$；形式 $b+a^2$ 限制到 $L$ 后在 $p$ 二阶消失，拉回为 $u(v+u)$。对(c)，$df_x$ 的切像也是 $V(b)$，$X$ 在 $x$ 光滑；取 $n=3$，形式 $b+a^2+a^3$ 的拉回为
\[
 u(v+u+u^2).
\]
这些显式式子说明了固定分支的机制；泛成员可约的结论来自整族均含 $E$ 且有非空移动部分，而不是仅凭三个特殊成员可约。原(a)--(c)的共同假设不能排除这个例子。9.2另行写出的非分歧条件不是9.3的默认假设。

\subsection*{三、9.3加强版的依据与边界}
正文新增
\[
 \dim f^{-1}(f(x))\le\dim X-2.
\]
这是一个易于使用的充分条件，不声称必要或恢复自手稿。三类线性族都包含所有在 $f(x)$ 至少二阶为零的 $n$ 次形式；其基点支集恰为 $f^{-1}(f(x))$。在代数闭基变换并取约化后，非零截面的每个分支余维为1；新增条件保证没有任何分支被困在基点集里。这与M19-8E中给出的稠密性补证相同。

(a)中两个比值就是 $t_1,t_2$；(b)中可先选在 $L$ 上为零但在像上非零的坐标，再从其余坐标中选取与它代数无关者。它们给出去掉基点后的映射至少二维的像。相应的一般截面在这个开部几何不可约；稠密性再将结论传给整个泛截面。

(c)还有一个原稿略写的归约。基变换到代数闭域，以 $W=\overline{f(X)}_{\mathrm{red}}$、$m=\dim W$ 记像闭包。取 $x$ 的光滑整邻域，$f$ 在稠密开集上落在 $W$ 的光滑部分，故切映射泛秩至多为 $m$。其 $(m+1)$ 阶子式在约化邻域上恒为零，故 $\dim L=\operatorname{rank}(df_x)\le m$。如果 $W\subset L$，就有
\[
 m\le\dim L\le m,
 \qquad W=L.
\]
因此 $L$ 上的独立坐标确实在像的函数域中代数无关。若 $m=2$，使用原三次式；若 $m\ge3$，使用英译所载的三个平方。平方比值的代数无关性在正特征下也成立，因为对两个代数无关元取正整数次幂不会降低超越次数。这完成正文说明所需的归约，不推出原稿已划去的“$f$ 在 $x$ 非分歧”。

本文所用的Bertini不可约性及基点条件可与
\href{https://stacks.math.columbia.edu/tag/0G4C}{Stacks Project，Tag 0G4C}核对；该处将固定余维1分支的排除明列为条件。上述爆破反例、切像论证以及与原稿的逐项对应为本轮独立工作。结论是几何不可约，不自动包括约化、几何整或光滑。

\subsection*{四、10.1的闭包与直线判据}
对于 $P^\vee$ 中可构造集 $Z$，“充分一般的几何直线避开 $Z$”可由其闭包讨论。但“存在一条直线避开 $Z$”并不能代替“存在一条直线避开 $\overline Z$”：若 $H\subset\mathbf P^2$ 是直线，$p\in H$，取 $Z=H\setminus\{p\}$，则过 $p$ 的另一条直线避开 $Z$，而 $\operatorname{codim}\overline Z=1$。

对闭集 $\overline Z$，余维至少2时，一般几何直线与其不交；余维为1时，射影维数定理使每条几何直线都与其相交。因此存在一条不交闭包的几何直线是正确判据。若 $k$ 无限，相应Grassmann簇的非空开集含 $k$-点，故可取定义在 $k$ 上的直线。有限域上没有这个自动的 $k$-有理点结论。正文同时明确了参数射影空间至少一维这一存在直线的前提。

\subsection*{五、10.2与10.3：已审范围内的局部证明}
法稿10.3为已划去的草稿，作者只写了“几乎平凡”和“未详细做过的练习”。以下另给译者证明，不把它接入原作者证明段。关于普通正则列与Koszul正则列的区别及局部诺特等价，核对了
\href{https://stacks.math.columbia.edu/tag/062D}{Tag 062D}和
\href{https://stacks.math.columbia.edu/tag/0638}{Tag 0638}；本轮不利用它们无根据地宣称任意非诺特范围已经审结。

首先，对任意拟凝聚 $G$ 与 $u:G\to\mathcal O_X$，令 $J=u(G)$。分次代数满射
\[
 \operatorname{Sym}_{\mathcal O_X}G
 \twoheadrightarrow \bigoplus_{n\ge0}J^n
\]
给出闭浸入 $\widetilde X\hookrightarrow\mathbf P_X(G)$。在爆破上，$J\mathcal O_{\widetilde X}$ 为可逆理想；典范商 $G_{\widetilde X}\twoheadrightarrow J\mathcal O_{\widetilde X}$ 的核在经 $u$ 映到 $\mathcal O_{\widetilde X}$ 后为零。因此 $K=u_P(H)$ 在 $\widetilde X$ 上消失，闭浸入经过 $V(K)$。这证明第一断言，无须给 $G$ 加有限秩。

其次，假设 $X$ 局部诺特、$G$ 局部自由有限秩 $r$，并按正文约定在 $T$ 附近由正则序列 $a_1,\ldots,a_r$ 描述 $u$。在 $\mathbf P^{r-1}_A$ 的 $U_1=D_+(T_1)$ 图上，写 $t_i=T_i/T_1$。$V(K)$ 的方程为
\[
 a_i-a_1t_i=0\quad(2\le i\le r),
 \qquad
 C=A[t_2,\ldots,t_r]/(a_i-a_1t_i)_{i=2}^r.
\]
在 $T$ 上方的局部环中，序列
\[
 a_1,\ a_2-a_1t_2,\ldots,a_r-a_1t_r
\]
正则：第一元为非零因子，模去它后是原正则列余项的平坦扩张。在诺特局部环中可换序，因而后 $r-1$ 元为正则列，并且 $a_1$ 在商环 $C$ 的相应局部环中为非零因子。在 $T$ 之外，若这些方程成立，则某个 $a_i$ 可逆迫使 $a_1$ 可逆；此时结论直接由消元得到。

一般有满射 $C\twoheadrightarrow A[J/a_1]$，其核由 $a_1$ 的幂所杀；可对照
\href{https://stacks.math.columbia.edu/tag/052P}{Tag 052P}的仿射爆破代数描述。上面的非零因子结论使这个核为零。每张标准图同样处理，粘合后得 $V(K)=\widetilde X$，并证明 $H\to\mathcal O_P$ 在其零点附近正则。$r=1$ 时 $H=0$、爆破为 $X$，与结论相容；秩按开闭部分处理；秩0时 $G=0$、$J=0$，其相对Proj和 $\mathbf P(0)$ 都为空，也与结论相容。

对10.2取 $r=2$，超平面束的关联方程为 $a_1\lambda+a_2\mu=0$。经 $[T_1:T_2]=[-\mu:\lambda]$ 认同后，它正是上面的 $a_2T_1-a_1T_2=0$。因此束与爆破在 $X$ 上同构。它在 $X\setminus T$ 上必为恒等；爆破的该开部概形稠密，且靶在 $X$ 上分离，故这个 $X$-同构唯一。此证明使用明确的正则性范围，不将任意零点理想的关联簇误认成爆破。

\subsection*{六、10.4已恢复的条件与仍存的依赖}
法59的不等号是 $\le1$，并非旧译猜测出来的可能修正。长度2正则列不可能存在于维数至多1的诺特局部环的极大理想中，所以正则中心只能为空。空中心时爆破为 $X$，这与正文对正则截面的约定一致。

法60的几何纤维底域为 $k'$，“至多一个分歧点”明确限定在每个 $s$ 上方。长度2的局部纤维为 $k'[t]/(t^2)$；这不是任意特征下温和分歧的断言。正文把“$P$ 中曲线”按闭子概形解释并写成射影曲线；若只取局部闭曲线，移去一个点会破坏纤维长度为固定 $d$ 的描述。

9.2仍引向未完整恢复的8.8，9.5仍是待写的残文。因此10.2的拟议几何应用和10.4的一般存在性不能因局部爆破计算正确就标为全部审结。该未决依赖已写入正文脚注及工程的未决事项。

\subsection*{七、执行记录与检验性质}
本轮看过法53--60整页及关键裁图，逐项将可读法文与英译36--41对照；没有以OCR输出替代扫描证据。固定分支、切映射及局部爆破证明属于数学推理；另用符号计算核对了三条爆破代换等式，计算只检查代数恒等式，不充当定理证明。TeX编译、单文件一致性和页面渲染结果记录在第十九版整合检查中；它们与数学完备性是不同事项。
\par\endgroup
